\documentclass{beamer}
\usepackage{graphicx}
\setlength{\parskip}{0 pt}
\setbeamertemplate{caption}{\centering\raggedright\scriptsize\insertcaption\par}
\title{ ClusterMining: extracting core structures of metallic nanoparticles from the atomic pair distribution function }
\subtitle{  }
\author{ Simon J. L. Billinge }
\date{\today}
\begin{document}

\begin{frame}
\titlepage
\end{frame}

\section{ introduction }


\begin{frame}

\frametitle{ What is local structure? }






Three ideas underpin the rest of the talk.


\begin{itemize}

  \item crystals are described by an average structure

  \item real materials deviate from that average

  \item those deviations control the properties

\end{itemize}


We will return to each of these.




\end{frame}


\begin{frame}

\frametitle{ The nanostructure problem }

\begin{columns}
\column{ 0.5\textwidth}
\begin{figure}
\begin{center}
\makebox[\columnwidth][c]{\includegraphics[width=0.9\dimexpr\columnwidth+2cm\relax, angle=0 ]{ example_nanoparticle }}

\caption{ A nanoparticle with a disordered surface. }

\end{center}
\end{figure}
\column{ 0.5\textwidth}


\begin{itemize}

  \item the surface is not the same as the core

\end{itemize}


\end{columns}







\end{frame}


\section{ Local structure of materials }


\begin{frame}

\frametitle{ What is local structure? }






Three ideas underpin the rest of the talk.


\begin{itemize}

  \item crystals are described by an average structure

  \item real materials deviate from that average

  \item those deviations control the properties

\end{itemize}


We will return to each of these.




\end{frame}


\begin{frame}

\frametitle{ Measuring a pair distribution function }



\begin{figure}
\begin{center}
\makebox[\textwidth][c]{\includegraphics[width=0.8\dimexpr\textwidth+2cm\relax, angle=0 ]{ example_diffractometer }}

\caption{ The diffractometer used to collect the data. }

\end{center}
\end{figure}





\end{frame}


\begin{frame}[shrink=10]

\frametitle{ From diffraction pattern to PDF }


\begin{columns}
\column{ 0.6\textwidth}

The transform is done in three steps.


\begin{itemize}

  \item correct and normalize the measured intensity

  \item Fourier transform to real space

  \item fit a structure model to the result

\end{itemize}


\column{ 0.4\textwidth}
\begin{figure}
\begin{center}
\makebox[\columnwidth][c]{\includegraphics[width=0.9\dimexpr\columnwidth+2cm\relax, angle=0 ]{ example_transform }}

\caption{ The Fourier transform of the total scattering structure function. }

\end{center}
\end{figure}
\end{columns}






\end{frame}


\end{document}
